\section{Dataset}
In our research, we introduce ``PredEx'', significantly differentiating itself from existing datasets in Legal Natural Language Processing (L-NLP), particularly in the context of the Indian judiciary. This dataset is designed to address the limitations of previous datasets, which primarily focused on prediction tasks and offered limited annotations for explanations.

\subsection{Dataset Compilation} In the Data Compilation process, we initially gathered a substantial corpus of about 20,000 court judgments randomly from the Supreme Court of India and various High Courts, utilizing the IndianKanoon website\footnote{\url{https://indiankanoon.org/}}, a legal search engine widely recognized for its comprehensive database of Indian legal documents. The corpus underwent a meticulous annotation process, where our team of legal experts focused on annotating explanations for the judgments. These annotations involved identifying and highlighting key sentences or segments within the case proceedings that significantly influenced the predicted outcomes, as well as providing reasoning for the judgments. Through this process, the original corpus was distilled to approximately 16,000 case files, each richly annotated with expert legal explanations.

Our approach in compiling the PredEx dataset was to randomly select cases, ensuring a broad representation across various types of judgments and legal decisions. This method was deliberately chosen to avoid bias towards any specific temporal aspect or domain. By adopting a random selection process, we aimed to capture the diverse nature of legal cases in the Indian judiciary system. This diversity was crucial to cover various aspects of law and legal decision-making, thereby enhancing the representativeness and applicability of our dataset for training AI models in legal judgment prediction and explanation.

Subsequent to the annotation phase, we undertook a preprocessing step to refine the dataset further. This preprocessing involved the removal of cases that were either too brief or where the final decision segments were challenging to discern. Such preprocessing is crucial for ensuring the quality and consistency of the data, particularly for training robust and reliable AI models; otherwise, it could introduce noise or bias into the model training. As a result of this preprocessing, the total number of case files in our dataset was reduced to 15,222 and is further divided into training and testing sets. We adopted an 80-20 split ratio for this purpose, ensuring a substantial volume of data for model training while still retaining a robust set for testing. Specifically, the training set consists of 12,178 documents, and the test set comprises 3,044 documents. 

In terms of balancing the test set, special attention was given to ensure fairness and representativeness in model evaluation. We carefully curated the test set to include a diverse range of case outcomes, such as different types of judgments and legal decisions. This diversity was not just in terms of the nature of cases but also in terms of the outcomes - for instance, balancing cases where appeals were accepted versus those that were dismissed. Such a balanced composition is crucial in avoiding biases towards any particular type of judgment and ensures that our AI models are tested against a wide spectrum of legal scenarios. This balanced nature of the test set is particularly important for maintaining the validity of our experiments and for ensuring the reliability and generalizability of our model's performance. These carefully processed and curated case files now form the core of our \texttt{PredEx} dataset, offering a rich resource for the Court Judgment Prediction and Explanation (CJPE) task. Detailed statistics of the final dataset, post-preprocessing, are presented in the following Table \ref{data-stats}.

%%%%%%%%%%%%%%%%%%%%%%%%%%%%%%%%%%%%%%%
\begin{table}[t]
  \centering
\resizebox{\columnwidth}{!}{%
  \begin{tabular}{lcc}
    \toprule
    & \textbf{Train} & \textbf{Test} \\
    \midrule
    No. of documents & 12,178 & 3,044 \\
    Average no. of tokens & 4,586 & 4,422 \\
    Minimum no. of tokens & 176 & 184 \\
    Maximum no. of tokens & 117,733 & 83,657 \\
    Acceptance percentage & 53.44\% & 50.00\% \\
    \bottomrule
  \end{tabular}
  }
    \caption{PredEx Statistics.}
  \label{data-stats}
\end{table}

%%%%%%%%%%%%%%%%%%%%%%%%%%%%%%%%%%%%%%%
\subsection{Annotation Process}
\subsubsection{Expert Involvement} We engaged a team of 10 legal experts, primarily law students in their 3rd and 4th years, from various Indian law colleges. These experts were selected based on their academic standing and understanding of legal processes, ensuring high-quality annotations.

\subsubsection{Annotation Timeline} The annotation process spanned from April 1, 2022, to October 30, 2023. This extensive period allowed for meticulous and thorough annotation, considering the complexity and detail required in legal document analysis.

\subsubsection{Work Allocation}
In our annotation process, each student was assigned around 30 judgment documents weekly, striking a balance between efficiency and the need for thorough, accurate annotations. This workload allocation enabled students to devote adequate time to each document, fostering precise and insightful annotations. 
\subsubsection{Role of Student Annotators}
The role of our student annotators was to meticulously identify and extract specific segments from the judgments that were pivotal to the judge's reasoning, rather than interpreting or analyzing these segments with their own legal reasoning. Their task was to pinpoint these key sections accurately, ensuring that the extracts faithfully represented the judicial reasoning as stated in the case documents. This extractive approach was critical to maintain the integrity and authenticity of the annotations, allowing the dataset to accurately reflect the content of the original legal texts without the introduction of subjective interpretations by the annotators.

% To ensure the robustness and reliability of these annotations, we implemented a systematic quality control process. Disagreements among annotators or uncertainties in annotations were addressed through a review mechanism overseen by a senior legal expert. This expert not only provided additional scrutiny to the annotations but also acted as a mediator to resolve any discrepancies. This process ensured a consistent and high-quality standard across all annotated documents. Regular training sessions and review meetings were also conducted to align the understanding and approach of all annotators, further enhancing the reliability of the dataset.
% 
\subsubsection{Annotation Quality Control Mechanism}
To guarantee the accuracy and consistency of the annotations, we implemented a comprehensive quality control system. Initially, each document was reviewed by a single annotator to ensure a consistent interpretation of the judicial content. Recognizing the complexity of legal texts and the potential for subjective interpretation, we established several layers of review to enhance the reliability of our annotations:
\begin{itemize}
    \item \textbf{Senior Expert Review:} Any disagreements or uncertainties in annotations were promptly escalated to a specialized review panel led by senior legal experts. These experts not only provided additional scrutiny but also mediated discrepancies among the initial annotations. Their extensive experience in legal practice and education enabled them to provide decisive and informed resolutions to any contentious or ambiguous annotations.
    \item \textbf{Regular Training and Meetings:} To further ensure consistency across annotations, regular training sessions and review meetings were conducted. These sessions served to align annotators on the legal framework and annotation criteria, reducing variability and enhancing the uniformity of the annotation process. Training included detailed discussions on identifying key legal arguments and the rationale within the judgments, which are critical for both the prediction and explanation aspects of our dataset.
\end{itemize}

This rigorous quality control process has ensured that our dataset meets high standards of reliability and validity. The annotations not only reflect the factual content of the legal decisions but also the detailed reasoning behind these judgments, making our dataset a robust resource for training and evaluating AI models in legal judgment prediction and explanation.

\subsubsection{Focus on Prediction and Explanations}
Diverging from previous datasets that primarily concentrate on the task of prediction, our \texttt{PredEx} dataset spans both prediction and explanations. The annotations in our dataset serve a dual purpose. Firstly, they identify the outcomes of the cases, fulfilling the prediction aspect. More importantly, they go a step further by providing detailed explanations behind these outcomes. These explanations elucidate the rationale or the legal reasoning that underpins the judgments. This dual emphasis on prediction and explanations fills a significant void in existing legal datasets. Typically, in other datasets, the aspect of explanation is either absent or not explored in depth. By contrast, \texttt{PredEx} enriches the field of legal AI with comprehensive annotations that shed light not just on what the judicial decisions are, but crucially, why these decisions were made. This focus on explanations is particularly vital, as it contributes to a more transparent and understandable AI-driven legal decision-making process.

\subsubsection{Largest Explainable Dataset} As a result of this extensive and detailed annotation process, we are releasing what is arguably the largest annotated dataset for legal judgment prediction and explanation in the Indian context. The size and comprehensiveness of this dataset set it apart from existing datasets in the field.

Our dataset represents a significant advancement in legal NLP, particularly for research and applications pertaining to the Indian judiciary. By providing a large-scale, richly annotated dataset that encompasses both prediction and explanation, we aim to facilitate more nuanced and sophisticated AI models capable of understanding and interpreting legal texts in a manner akin to human legal experts. This dataset is not only a resource for advancing AI technology in the legal domain but also a step towards enhancing transparency and accountability in AI-assisted legal decision-making. 
